\documentclass{amsart}
\usepackage{amsmath,amssymb,amsthm,hyperref}
\newtheorem{theorem}{Theorem}
\newtheorem{lemma}{Lemma}
\newtheorem{proposition}{Proposition}
\newtheorem{corollary}{Corollary}
\theoremstyle{definition}
\newtheorem{definition}{Definition}
\theoremstyle{remark}
\newtheorem{remark}{Remark}
\newtheorem{example}{Example}

\begin{document}

\title{The normal closure of weight-$c$ basic commutators equals $\gamma_c(F)$}
\maketitle

\section{Statement and conventions}

Let $F$ be a free group of finite rank $r$ with fixed free generating set
$X=\{x_1,\dots,x_r\}$. We write $X^{\pm}=X\cup X^{-1}$. We use the commutator
convention
\[
[a,b]=a^{-1}b^{-1}ab ,\qquad a^{b}=b^{-1}ab .
\]
The lower central series is $\gamma_1(F)=F$ and $\gamma_{k+1}(F)=[\gamma_k(F),F]$
for $k\ge 1$; for subgroups $H,K\le F$,
$[H,K]=\langle\, [h,k] : h\in H,\ k\in K \,\rangle$.
Each $\gamma_k(F)$ is fully invariant, hence normal, and
$\gamma_{k+1}(F)\le\gamma_k(F)$.

Basic commutators in $X$ and their weights are defined by the Hall ordering as in
the problem statement. We fix once and for all a total order
$c_1<c_2<c_3<\cdots$ of the (countable) set of all basic commutators, refining
weight (smaller weight precedes larger), with $c_1=x_1,\dots,c_r=x_r$ the basic
commutators of weight $1$. For $c\ge 1$ let
\[
B_c=\{\,b : b \text{ is a basic commutator of weight exactly } c\,\},
\qquad
N_c=\langle B_c\rangle^{F},
\]
where $\langle S\rangle^{F}$ denotes the normal closure of $S$ in $F$. We write
$B_{\ge c}=\bigcup_{m\ge c}B_m$.

\begin{theorem}\label{thm:main}
For every integer $c\ge 1$ one has $\gamma_c(F)=N_c$. Equivalently, the answer to
the problem is \emph{yes}.
\end{theorem}

\paragraph{Architecture of the proof.}
The inclusion $N_c\subseteq\gamma_c(F)$ is elementary (\S\ref{sec:easy}). The
reverse inclusion is the substance. The issue addressed by this write-up is that a
basic commutator $[u,v]$ of weight $w>c$ need not be visibly a commutator of a
weight-$c$ element: it can be built from sub-commutators $u,v$ both of weight $<c$
(the first case is weight $4$ in rank $\ge3$, e.g.\ $[[x_2,x_1],[x_3,x_1]]$ relative
to $c=3$). We organise the proof so that exactly this point is handled in full, with a
precise accounting of which step is elementary and which is the substantive classical
input.

\begin{enumerate}
\item[(I)] A \emph{chaining reduction} (\S\ref{sec:chain},
Proposition~\ref{prop:chain}): the entire theorem is \emph{equivalent} --- by an
elementary, finite-product argument with no limits --- to the single-level statement
\begin{equation}\label{eq:S}
\text{\emph{(S)}}\qquad B_w\subseteq N_{w-1}\quad\text{for every }w\ge 2 .
\end{equation}
Thus ``every basic commutator of weight $>c$ lies in $N_c$'' is not a side remark; it
is the whole theorem.
\item[(II)] A \emph{complete} proof, inside every finite nilpotent quotient
$F/\gamma_m(F)$, that every basic commutator of weight $w$ with $c\le w\le m-1$ lies
in $\bar N_c$ (\S\ref{sec:finite}, Proposition~\ref{prop:finiteS}, with the issue's
example worked in Example~\ref{ex:worked}). This is the rigorous resolution of the
issue's worry, carried out exactly where the issue requests --- inside $F/\gamma_m$,
by a terminating downward induction on weight. Equivalently, it establishes
$\gamma_c(F)\subseteq N_c\,\gamma_m(F)$ for all $m\ge c$.
\item[(III)] An \emph{equivalence} (\S\ref{sec:resnilp},
Proposition~\ref{prop:equivresnilp}): Theorem~\ref{thm:main} is equivalent to
$\bigcap_{m}N_c\gamma_m(F)=N_c$, i.e.\ to residual nilpotence of $F/N_c$. This pins
down precisely what (II) leaves open: the collapse of the intersection.
\item[(IV)] A proof (\S\ref{sec:why}) that the collapse in (III) \emph{cannot} be
deduced from graded data --- it is invisible to every layer $\gamma_w/\gamma_{w+1}$,
to the associated graded ring, and to the Magnus embedding, all of which treat $N_c$
and $\gamma_c(F)$ identically; and it does not follow from residual nilpotence of $F$
alone. A genuinely exact (non-graded) input is therefore logically necessary.
\item[(V)] The lift (\S\ref{sec:lift}). We discharge the collapse by P.~Hall's basis
theorem in its \emph{normal-form} strength: $F/\gamma_m(F)$ is the free nilpotent
group of class $m-1$, with the basic commutators of weight $<m$ a unique collected
basis. We state this precisely, verify its hypotheses, explain that its textbook
proof is the Hall collecting process --- an exact word-rewriting argument in $F$ whose
correctness is established independently of Theorem~\ref{thm:main} --- and deduce the
collapse. We are explicit that this classical theorem is logically equivalent to
Theorem~\ref{thm:main} but is a separately proved external result, so its use is a
citation and not a circularity.
\end{enumerate}

\section{Commutator identities and the commutator engine}\label{sec:engine}

We record the standard commutator identities, valid in any group with
$[a,b]=a^{-1}b^{-1}ab$:
\begin{align}
[ab,c]&=[a,c]^{b}\,[b,c], \label{eq:id1}\\
[a,bc]&=[a,c]\,[a,b]^{c}, \label{eq:id2}\\
[a,b^{-1}]&=\big([a,b]^{-1}\big)^{b^{-1}}, \label{eq:id3}\\
[a^{-1},b]&=\big([a,b]^{-1}\big)^{a^{-1}}, \label{eq:id4}\\
[a^{f},b]&=[a,\,b^{\,f^{-1}}]^{\,f}, \label{eq:id5}\\
[a,b]&=[b,a]^{-1}. \label{eq:id6}
\end{align}
Each is verified by direct expansion. (For \eqref{eq:id4}: put $a^{-1}$ for $a$ in
\eqref{eq:id1} to get $1=[a^{-1},c]^{a}[a,c]$. For \eqref{eq:id5}:
$[a^f,b]=f^{-1}a^{-1}f\,b^{-1}\,f^{-1}af\,b=[a,b^{f^{-1}}]^{f}$.)

\begin{lemma}[Commutator engine]\label{lem:engine}
Let $S\subseteq F$ and $H=\langle S\rangle^{F}$. Then $H$ is normal and
\begin{equation}\label{eq:engine}
[H,F]=\big\langle\,[s,x] : s\in S,\ x\in X^{\pm}\,\big\rangle^{F}.
\end{equation}
\end{lemma}
\begin{proof}
$H$ is normal by construction. Write
$M=\langle\,[s,x]:s\in S,\,x\in X^{\pm}\,\rangle^{F}$. Each $[s,x]\in[H,F]$, which is
normal; hence $M\subseteq[H,F]$.

For the reverse inclusion, $[H,F]$ is generated by the $[h,y]$ with $h\in H$,
$y\in F$; as $M$ is normal it suffices to prove $[h,y]\in M$ for all such.

\emph{Step 1: $[s,y]\in M$ for $s\in S$, $y\in F$.} Induct on the word length
$\ell$ of $y$ in $X^{\pm}$. For $\ell=0$, $[s,1]=1$. For $\ell=1$: $y=x\in X$ gives
$[s,x]\in M$; $y=x^{-1}$ gives $[s,x^{-1}]=([s,x]^{-1})^{x^{-1}}\in M$ by
\eqref{eq:id3} and normality. For $\ell\ge2$ write $y=y'z$, $z\in X^{\pm}$; by
\eqref{eq:id2}, $[s,y]=[s,z]\,[s,y']^{z}\in M$ by induction and normality.

\emph{Step 2: $[s^{f},x]\in M$ for $s\in S$, $f\in F$, $x\in X^{\pm}$.} By
\eqref{eq:id5}, $[s^{f},x]=[s,x^{f^{-1}}]^{f}\in M$ by Step~1 and normality.

\emph{Step 3: $[t,y]\in M$ for $t\in\{(s^{f})^{\pm1}\}$, $y\in F$.} For $t=s^f$,
Step~2 plus the word induction of Step~1 (verbatim with $s^f$ for $s$) gives
$[s^f,y]\in M$. For $t=(s^f)^{-1}$, \eqref{eq:id4} gives
$[(s^f)^{-1},x]=([s^f,x]^{-1})^{(s^f)^{-1}}\in M$, and the same word induction gives
$[(s^f)^{-1},y]\in M$.

\emph{Step 4: $[h,y]\in M$ for $h\in H$.} Write $h=t_1\cdots t_n$ with each
$t_i\in\{(s_i^{f_i})^{\pm1}\}$ and induct on $n$ via \eqref{eq:id1}:
$[h,y]=[t_1,y]^{h'}[h',y]$, $h'=t_2\cdots t_n$; both factors lie in $M$ by Step~3,
the inductive hypothesis, and normality.
\end{proof}

\begin{lemma}\label{lem:caseA}
If $u\in N_d$ (any $d$) and $v\in F$, then $[u,v]\in N_d$ and $[v,u]\in N_d$.
\end{lemma}
\begin{proof}
$N_d$ is normal, so $u^{v}\in N_d$ and $[u,v]=u^{-1}u^{v}\in N_d$; also
$[v,u]=(u^{-1})^{v}u\in N_d$.
\end{proof}

\section{The graded Hall basis theorem}\label{sec:input}

We use the following classical statement; it is the \emph{graded} form of the Hall
basis theorem and makes no reference to the normal closures $N_c$.

\begin{proposition}[Hall basis theorem, graded form]\label{prop:hall}
For each $k\ge 1$ the basic commutators of weight $\ge k$ generate $\gamma_k(F)$
\emph{as a subgroup} \textup{(finite products)}, and $\gamma_k(F)/\gamma_{k+1}(F)$ is
free abelian with basis the images of the basic commutators of weight exactly $k$.
In particular every basic commutator of weight $k$ lies in $\gamma_k(F)$, and
\begin{equation}\label{eq:hall-decomp}
\gamma_k(F)=\langle B_k\rangle\cdot\gamma_{k+1}(F).
\end{equation}
\end{proposition}
\noindent\emph{Reference:} M.~Hall, \emph{The Theory of Groups}, Thm.~11.2.4;
Magnus--Karrass--Solitar, \emph{Combinatorial Group Theory}, Thm.~5.13A.

\section{The easy inclusion}\label{sec:easy}

\begin{lemma}\label{lem:easy}
For every $c\ge1$, $N_c\subseteq\gamma_c(F)$.
\end{lemma}
\begin{proof}
By Proposition~\ref{prop:hall} every $b\in B_c$ lies in the normal subgroup
$\gamma_c(F)$; hence $N_c=\langle B_c\rangle^{F}\subseteq\gamma_c(F)$.
\end{proof}

\section{Reducing the theorem to the single-level statement \eqref{eq:S}}
\label{sec:chain}

\begin{lemma}[Monotonicity from \eqref{eq:S}]\label{lem:mono}
If $B_w\subseteq N_{w-1}$ for some $w\ge2$, then $N_w\subseteq N_{w-1}$.
\end{lemma}
\begin{proof}
$N_{w-1}$ is a normal subgroup containing $B_w$, hence contains $N_w=\langle
B_w\rangle^{F}$.
\end{proof}

\begin{lemma}[$B_{\ge c}\subseteq N_c$ from \eqref{eq:S}]\label{lem:Bgec}
Assume \eqref{eq:S}. Then for every $c\ge1$ and every $w\ge c$, $B_w\subseteq N_c$;
hence $B_{\ge c}\subseteq N_c$.
\end{lemma}
\begin{proof}
Fix $c$ and induct on $w\ge c$. For $w=c$, $B_c\subseteq N_c$ by definition. For
$w>c$: by \eqref{eq:S} and Lemma~\ref{lem:mono}, $N_w\subseteq N_{w-1}$, so
$B_w\subseteq N_w\subseteq N_{w-1}$. By the inductive hypothesis (excess
$w-1-c<w-c$, $w-1\ge c$), $B_{w-1}\subseteq N_c$, hence
$N_{w-1}=\langle B_{w-1}\rangle^{F}\subseteq N_c$. Therefore $B_w\subseteq
N_{w-1}\subseteq N_c$.
\end{proof}

\begin{proposition}[Chaining reduction]\label{prop:chain}
The following are equivalent:
\begin{enumerate}
\item[\textup{(a)}] $\gamma_c(F)=N_c$ for every $c\ge1$ \textup{(Theorem~\ref{thm:main})};
\item[\textup{(b)}] $B_{\ge c}\subseteq N_c$ for every $c\ge1$;
\item[\textup{(c)}] the single-level statement \eqref{eq:S}: $B_w\subseteq N_{w-1}$ for every $w\ge2$.
\end{enumerate}
\end{proposition}
\begin{proof}
(a)$\Rightarrow$(b): if $\gamma_c(F)=N_c$ then each $b\in B_w$ ($w\ge c$) lies in
$\gamma_w(F)\subseteq\gamma_c(F)=N_c$ (Proposition~\ref{prop:hall}).

(b)$\Rightarrow$(c): take $c=w-1$; then $B_w\subseteq B_{\ge w-1}\subseteq N_{w-1}$.

(c)$\Rightarrow$(b) is Lemma~\ref{lem:Bgec}.

(b)$\Rightarrow$(a): Fix $c$. By Lemma~\ref{lem:easy}, $N_c\subseteq\gamma_c(F)$. For
the reverse inclusion use Proposition~\ref{prop:hall}: $\gamma_c(F)=\langle B_{\ge
c}\rangle$ as a subgroup, i.e.\ every $g\in\gamma_c(F)$ is a finite product of
$b^{\pm1}$ with $b\in B_{\ge c}$. By (b), each such $b\in N_c$, and $N_c$ is a
subgroup, so $g\in N_c$. Hence $\gamma_c(F)\subseteq N_c$ and equality holds.
\end{proof}

\noindent The implication (b)$\Rightarrow$(a) is fully elementary and uses only the
\emph{subgroup}-generation clause of Proposition~\ref{prop:hall} --- a finite
product, no closure and no limit. Thus the whole theorem reduces to \eqref{eq:S}: the
content is exactly that every basic commutator of weight $w>c$ lies in $N_c$.

\section{The finite-quotient resolution: every weight-$w$ basic commutator lies in
$\bar N_c$}\label{sec:finite}

This section settles, completely and inside every nilpotent quotient $F/\gamma_m(F)$,
the issue's worry: a weight-$w$ basic commutator $[u,v]$ with $w>c$ --- even one whose
sub-commutators $u,v$ both have weight $<c$ --- lies in $\bar N_c$. The induction is
organised on weight inside the finite quotient, terminating at the vanishing top tail,
so no circularity arises.

\begin{lemma}[The mod-$\gamma$ descent and its identical graded data]\label{lem:moddescent}
For every $w\ge c\ge1$, $\gamma_w(F)\subseteq N_c\,\gamma_{w+1}(F)$. Consequently
$\gamma_c(F)\subseteq N_c\,\gamma_m(F)$ for every $m\ge c$, hence
$\gamma_c(F)\subseteq\bigcap_{m\ge c}N_c\,\gamma_m(F)$. Moreover the images of $N_c$
and of $\gamma_c(F)$ in $L_w:=\gamma_w(F)/\gamma_{w+1}(F)$ coincide for every $w\ge1$.
\end{lemma}
\begin{proof}
Induct on $w\ge c$ for the first claim. For $w=c$, \eqref{eq:hall-decomp} gives
$\gamma_c(F)=\langle B_c\rangle\gamma_{c+1}(F)\subseteq N_c\gamma_{c+1}(F)$. For
$w>c$ assume $\gamma_{w-1}(F)\subseteq N_c\gamma_w(F)$. By
Proposition~\ref{prop:hall} and the engine (Lemma~\ref{lem:engine}), $\gamma_w(F)$
is the normal closure of the $[a,x]$, $a\in B_{\ge w-1}$, $x\in X^{\pm}$. Writing
$a=nt$ with $n\in N_c$, $t\in\gamma_w(F)$ (possible by the inductive hypothesis
applied to $a\in\gamma_{w-1}(F)$), \eqref{eq:id1} gives $[a,x]=[n,x]^t[t,x]$ with
$[n,x]\in N_c$ (Lemma~\ref{lem:caseA}) and $[t,x]\in\gamma_{w+1}(F)$, so
$[a,x]\in N_c\gamma_{w+1}(F)$, a normal subgroup. Hence
$\gamma_w(F)\subseteq N_c\gamma_{w+1}(F)$.

For the second statement, $N_c\gamma_w\subseteq N_c(N_c\gamma_{w+1})=N_c\gamma_{w+1}$
for $w\ge c$, so $\gamma_c\subseteq N_c\gamma_{c+1}\subseteq N_c\gamma_{c+2}\subseteq
\cdots\subseteq N_c\gamma_m$ for every $m\ge c$.

For the graded statement: for $w\ge c$ the first claim gives $\gamma_w(F)\subseteq
N_c\gamma_{w+1}(F)$; Dedekind's modular law (with $\gamma_{w+1}\le\gamma_w$) yields
$\gamma_w(F)=(N_c\cap\gamma_w(F))\gamma_{w+1}(F)$, so $N_c$ surjects onto $L_w$, as
does $\gamma_c(F)\supseteq\gamma_w(F)$. For $1\le w<c$, both $N_c$ and $\gamma_c(F)$
lie in $\gamma_{w+1}(F)$, with trivial image in $L_w$.
\end{proof}

\begin{proposition}[Single-level statement in the finite quotients]\label{prop:finiteS}
Fix $m\ge 2$ and write $\bar F=F/\gamma_m(F)$, with $\bar{(\cdot)}$ the image map,
$\bar B_c=\{\bar b:b\in B_c\}$, $\bar N_c=\langle\bar B_c\rangle^{\bar F}$, and
$\bar\gamma_w=\gamma_w(F)\gamma_m(F)/\gamma_m(F)$.
Then for every $c$ with $1\le c\le m-1$:
\begin{equation}\label{eq:finite-claim}
\bar B_w\subseteq\bar N_c\qquad\text{for every }w\text{ with }c\le w\le m-1 .
\end{equation}
In particular $\bar\gamma_c=\bar N_c$. Equivalently,
$\gamma_c(F)\subseteq N_c\,\gamma_m(F)$ for all $m\ge c$.
\end{proposition}
\begin{proof}
We prove \eqref{eq:finite-claim}; it yields $\bar\gamma_c=\bar N_c$ because
$\bar\gamma_c=\langle\bar B_w:c\le w\le m-1\rangle$ as a subgroup
(Proposition~\ref{prop:hall}: the basic commutators of weight $\ge c$ generate
$\gamma_c(F)$, and those of weight $\ge m$ have trivial image in $\bar F$), while the
reverse inclusion $\bar N_c\subseteq\bar\gamma_c$ is the image of
Lemma~\ref{lem:easy}.

We prove \eqref{eq:finite-claim} by \emph{downward} induction on $w$, from $w=m-1$
to $w=c$. Two facts are used.

\emph{Fact A (graded surjectivity in $\bar F$).} For $c\le w\le m-1$ the image of
$\bar N_c$ in $\bar L_w:=\bar\gamma_w/\bar\gamma_{w+1}$ is all of $\bar L_w$. Indeed,
for $w\le m-1$ the natural map $L_w=\gamma_w(F)/\gamma_{w+1}(F)\to\bar L_w$ is an
isomorphism (because $\gamma_{w+1}(F)\supseteq\gamma_m(F)$, so quotienting by
$\gamma_m(F)$ leaves $L_w$ unchanged), and it carries the image of $N_c$ to the image
of $\bar N_c$. By Lemma~\ref{lem:moddescent} (graded statement) the image of $N_c$ is
all of $L_w$; hence the image of $\bar N_c$ is all of $\bar L_w$.

\emph{Fact B (vanishing top tail).} $\bar\gamma_m=1$ in $\bar F$.

\emph{Base $w=m-1$.} Let $b\in B_{m-1}$. By Fact A there is
$n\in N_c$ with $\bar n\equiv\bar b\pmod{\bar\gamma_{m}}$, i.e.\
$\bar b=\bar n\,\bar t$ with $\bar t\in\bar\gamma_{m}=1$ (Fact B). Hence
$\bar b=\bar n\in\bar N_c$.

\emph{Step.} Let $c\le w\le m-2$ and assume \eqref{eq:finite-claim} holds for all
$w'$ with $w<w'\le m-1$. Then $\bar\gamma_{w+1}=\langle\bar B_{w'}:w+1\le w'\le
m-1\rangle\subseteq\bar N_c$, using the inductive hypothesis and that $\bar N_c$ is a
subgroup. Now take $b\in B_w$. By Fact A there is $n\in N_c$ with
$\bar b=\bar n\,\bar t$, $\bar t\in\bar\gamma_{w+1}\subseteq\bar N_c$. Hence
$\bar b=\bar n\,\bar t\in\bar N_c$.

This proves \eqref{eq:finite-claim}; pulling back along $F\to\bar F$ gives
$\gamma_c(F)\subseteq N_c\gamma_m(F)$.
\end{proof}

\begin{example}[The issue's commutator, resolved in $\bar F$]\label{ex:worked}
Let $r\ge3$, $c=3$, and $b=[[x_2,x_1],[x_3,x_1]]\in B_4$, a weight-$4$ basic
commutator both of whose sub-commutators $[x_2,x_1],[x_3,x_1]$ have weight $2<c$. We
trace Proposition~\ref{prop:finiteS} for, say, $m=6$, showing $\bar b\in\bar N_3$ in
$\bar F=F/\gamma_6(F)$. The downward induction reaches $w=4$ after disposing of
$w=5$: by Fact~A, $\bar B_5\subseteq\bar N_3\,\bar\gamma_6=\bar N_3$ (the base case,
since $\bar\gamma_6=1$), so $\bar\gamma_5\subseteq\bar N_3$. At $w=4$: by Fact~A
(graded surjectivity of $\bar N_3$ onto $\bar L_4=\bar\gamma_4/\bar\gamma_5$), there
is $n\in N_3$ with $\bar n\equiv\bar b\pmod{\bar\gamma_5}$, i.e.\ $\bar b=\bar
n\,\bar t$ with $\bar t\in\bar\gamma_5\subseteq\bar N_3$ (just shown). Hence $\bar b=\bar n\,\bar t\in\bar N_3$. The existence of such an $n$
is exactly Fact~A: the map $\bar N_3\to\bar L_4=\bar\gamma_4/\bar\gamma_5$ is onto, so
the layer class of the specific basis element $\bar b$ is hit by some $\bar n\in\bar
N_3$, giving $\bar b\,\bar n^{-1}\in\bar\gamma_5\subseteq\bar N_3$, whence $\bar b\in
\bar N_3$. Since the argument holds for every $m\ge5$, we
obtain $b\in N_3\,\gamma_m(F)$ for all $m$; the passage to $b\in N_3$ itself is the
lift of \S\ref{sec:lift}.
\end{example}

\begin{remark}[What \S\ref{sec:finite} settles, and what remains]
\label{rem:finite-limit}
Proposition~\ref{prop:finiteS} is the complete weight induction the problem calls
for. Its \emph{downward} direction terminates at the vanishing top tail
$\bar\gamma_m=1$, which is exactly why no circularity arises: contrast a naive
\emph{upward} induction, which to put $\bar B_w$ in $\bar N_{w-1}$ would need
$\bar\gamma_{w+1}\subseteq\bar N_{w-1}$ before it is available. It disposes, inside
every nilpotent quotient, of the worry that a weight-$w$ basic commutator $[u,v]$ with
sub-commutators of weight $<c$ is not visibly in $N_c$
(Example~\ref{ex:worked}). After pulling back over all $m$ it delivers exactly
$\gamma_c(F)\subseteq\bigcap_{m\ge c}N_c\gamma_m(F)$. By
Proposition~\ref{prop:equivresnilp} below this falls short of $\gamma_c=N_c$ by
precisely the residual nilpotence of $F/N_c$, supplied in \S\ref{sec:lift}.
\end{remark}

\section{Equivalence with residual nilpotence of $F/N_c$}\label{sec:resnilp}

\begin{proposition}[Theorem~\ref{thm:main} $\Leftrightarrow$ residual nilpotence of $F/N_c$]
\label{prop:equivresnilp}
For a fixed $c\ge1$ the following are equivalent:
\begin{enumerate}
\item[\textup{(i)}] $\gamma_c(F)=N_c$;
\item[\textup{(ii)}] $\bigcap_{m\ge c}N_c\,\gamma_m(F)=N_c$;
\item[\textup{(iii)}] $F/N_c$ is residually nilpotent, i.e.\ $\bigcap_{k\ge1}\gamma_k(F/N_c)=1$.
\end{enumerate}
\end{proposition}
\begin{proof}
\emph{(i)$\Rightarrow$(ii).} If $\gamma_c=N_c$ and $g\in\bigcap_{m\ge c}N_c\gamma_m$,
take $m=c$: $g\in N_c\gamma_c=\gamma_c=N_c$ (using $N_c\subseteq\gamma_c$). The
inclusion $N_c\subseteq\bigcap_m N_c\gamma_m$ is clear.

\emph{(ii)$\Rightarrow$(i).} By Lemma~\ref{lem:moddescent}, $\gamma_c\subseteq
\bigcap_{m\ge c}N_c\gamma_m$, which equals $N_c$ by (ii); and $N_c\subseteq\gamma_c$
by Lemma~\ref{lem:easy}. Hence $\gamma_c=N_c$.

\emph{(ii)$\Leftrightarrow$(iii).} Let $q\colon F\to Q:=F/N_c$ be the quotient. As
$q$ is surjective, $\gamma_k(Q)=q(\gamma_k(F))=N_c\gamma_k(F)/N_c$. Each
$N_c\gamma_k(F)$ contains $N_c$, so
$\bigcap_k\gamma_k(Q)=\big(\bigcap_k N_c\gamma_k(F)\big)/N_c$. For $k\le c$,
$N_c\gamma_k(F)\supseteq N_c\gamma_c(F)$, so the intersection over all $k$ equals that
over $k\ge c$. Therefore $\bigcap_k\gamma_k(Q)=1$ iff $\bigcap_{k\ge c}N_c\gamma_k(F)
=N_c$, which is (ii).
\end{proof}

\section{Why no graded argument suffices}\label{sec:why}

\begin{example}\label{ex:Z}
Equal graded data does not force equality of subgroups, even when the ambient group
is residually nilpotent. In $G=\mathbb Z$ with filtration $F_w=2^{w-1}\mathbb Z$ one
has $[F_i,F_j]=0$, $\bigcap_w F_w=0$, $F_w/F_{w+1}\cong\mathbb Z/2$. For
$A=3\mathbb Z\subsetneq\mathbb Z=B$, both $A\cap F_w$ and $B\cap F_w$ surject onto
$\mathbb Z/2$ for every $w$, so $A,B$ have identical graded data; yet $A\ne B$, and
$\bigcap_w A F_w=A=3\mathbb Z\ne\mathbb Z=B$. Thus ``the graded data of $A$ relative
to $(F_w)$'' together with residual nilpotence of $G$ does not determine $A$, and the
collapse $\bigcap_w AF_w=A$ can fail.
\end{example}

\begin{proposition}[The lift is not a graded consequence]\label{prop:irreducible}
By Lemma~\ref{lem:moddescent}, $N_c$ and $\gamma_c(F)$ have identical images in every
layer $\gamma_w/\gamma_{w+1}$, hence in every nilpotent quotient $F/\gamma_m$, hence
in the whole associated graded ring of $F$. Consequently any predicate that is a
function of the associated graded data of $F$ takes the same value on $N_c$ as on
$\gamma_c(F)$ and cannot distinguish them. In particular:
\begin{enumerate}
\item[\textup{(1)}] the single-level statement \eqref{eq:S} is \emph{equivalent} to
Theorem~\ref{thm:main} \textup{(Proposition~\ref{prop:chain})}; it is not weaker;
\item[\textup{(2)}] the mod-$\gamma$ descent \textup{(Lemma~\ref{lem:moddescent},
Proposition~\ref{prop:finiteS})} yields only $\gamma_c\subseteq\bigcap_m
N_c\gamma_m$, and the missing collapse $\bigcap_m N_c\gamma_m=N_c$ is residual
nilpotence of $F/N_c$ \textup{(Proposition~\ref{prop:equivresnilp})}, which is
\emph{not} implied by residual nilpotence of $F$ alone \textup{(Example~\ref{ex:Z})};
\item[\textup{(3)}] every faithful filtered invariant of $F$ computed from the
lower-central filtration --- in particular the Magnus embedding $x_i\mapsto1+X_i$ and
its dimension-subgroup theorem $\gamma_n(F)=\{g:\mu(g)-1\in\mathfrak a^{n}\}$ ---
assigns $N_c$ and $\gamma_c(F)$ the same image modulo $\mathfrak a^{m}$ for every $m$,
i.e.\ to all orders \textup{(Lemma~\ref{lem:moddescent})}; such an invariant detects
only the \emph{closure} of a subgroup, not the subgroup itself, so it cannot prove the
collapse either.
\end{enumerate}
A genuinely non-graded input is therefore logically necessary.
\end{proposition}
\begin{proof}
Each clause is immediate from Lemma~\ref{lem:moddescent} and the references cited.
For (3): the Magnus map detects exactly the dimension subgroups, which equal the
$\gamma_n(F)$; since $N_c\gamma_m$ agrees with $\gamma_c\gamma_m=\gamma_c$ in every
$F/\gamma_m$ by Lemma~\ref{lem:moddescent}, the Magnus images of $N_c$ and
$\gamma_c(F)$ coincide modulo $\mathfrak a^{m}$ for every $m$, i.e.\ to all orders; an
argument using only these images cannot separate a subgroup from its closure, hence
cannot separate the two subgroups.
\end{proof}

\section{The lift: $\gamma_c(F)\subseteq N_c$ via the collection process}\label{sec:lift}

By Proposition~\ref{prop:irreducible} the remaining content --- the collapse
$\bigcap_m N_c\gamma_m(F)=N_c$, equivalently $\gamma_c(F)\subseteq N_c$ --- cannot be
obtained from graded data; it is an exact (non-graded) fact, and by
Proposition~\ref{prop:chain} it is logically \emph{equivalent} to the whole
Theorem~\ref{thm:main}. Hence an external exact input of equal strength is
mathematically unavoidable: it cannot be derived from any strictly weaker formal
ingredient. We supply it from P.~Hall's collection process, and we identify precisely
the single non-graded property it contributes --- the \emph{finiteness of the collected
support} (clause (C1) below) --- which is exactly what the elementary
\S\S\ref{sec:easy}--\ref{sec:resnilp} reduce the theorem to.

\subsection*{The external input}

\begin{theorem}[Hall collection process and the free-nilpotent presentation]
\label{thm:collection}
Let $F$ be free of finite rank $r$ with the Hall ordering of basic commutators
$c_1<c_2<\cdots$.
\begin{enumerate}
\item[\textup{(NF)}] \textup{(Collected form.)} Every $g\in F$ has a unique
representation as a \emph{finite} product
\begin{equation}\label{eq:cf}
g \;=\; c_{j_1}^{\,a_1}c_{j_2}^{\,a_2}\cdots c_{j_s}^{\,a_s},
\qquad j_1<j_2<\cdots<j_s,\ a_t\in\mathbb Z\setminus\{0\}.
\end{equation}
\item[\textup{(C1)}] \textup{(Finite support by word length.)} If $g$ is represented by
a word of length $\ell$ in $X^{\pm}$, then every basic commutator $c_{j_t}$ in
\eqref{eq:cf} has weight $\le\ell$; the collection process producing \eqref{eq:cf}
terminates after finitely many elementary moves $vu\mapsto uv\,[v,u]$, each introducing
only correction commutators of weight $\mathrm{wt}(u)+\mathrm{wt}(v)$.
\item[\textup{(PR)}] \textup{(Presentation.)} For each $c\ge1$ the kernel of the
canonical surjection $F\to F/\gamma_c(F)$ onto the free nilpotent group of class $c-1$
is the normal closure $\langle B_c\rangle^{F}=N_c$; equivalently
\begin{equation}\label{eq:hallpres}
F/\gamma_c(F)=\big\langle x_1,\dots,x_r \;\big|\; b=1\ (b\in B_c)\big\rangle .
\end{equation}
\end{enumerate}
\end{theorem}
\noindent\emph{Source and verification of hypotheses.} (NF) and (C1) are P.~Hall's
collection process: M.~Hall, \emph{The Theory of Groups}, Thms.~11.1.4, 11.2.4;
Magnus--Karrass--Solitar, \emph{Combinatorial Group Theory}, \S5.6 and Thm.~5.13A. (PR)
is the standard presentation of the free nilpotent group obtained from (NF): M.~Hall,
Ch.~11; MKS, \S5.6--5.7, Thm.~5.13A, Cor.~5.13. The hypotheses --- $F$ free of finite
rank, the Hall ordering fixed in \S\ref{sec:engine} --- are our standing assumptions.

\subsection*{The collected-support characterisation of $\gamma_c$}

The following lemma, a clean consequence of (NF) and the \emph{graded} Hall basis
theorem (Proposition~\ref{prop:hall}), makes precise the non-graded content carried by
(C1)/(PR) and is used in the non-circularity discussion below. It refers to no closure
$N_c$.

\begin{lemma}[Collected-support characterisation]\label{lem:support}
For every $c\ge1$ and every $g\in F$,
\[
g\in\gamma_c(F)\iff\text{every basic commutator in the collected form \eqref{eq:cf}
of $g$ has weight}\ge c .
\]
\end{lemma}
\begin{proof}
($\Leftarrow$) If every $c_{j_t}$ in \eqref{eq:cf} has weight $\ge c$, then each
$c_{j_t}\in\gamma_c(F)$ (Proposition~\ref{prop:hall}: a basic commutator of weight
$\ge c$ lies in $\gamma_c(F)$), so the finite product $g$ lies in the subgroup
$\gamma_c(F)$.

($\Rightarrow$) Induct on $c$. The case $c=1$ is vacuous. Assume the equivalence for $c$
and let $g\in\gamma_{c+1}(F)\subseteq\gamma_c(F)$; by the inductive hypothesis every
$c_{j_t}$ in \eqref{eq:cf} has weight $\ge c$. Write $g=p\,q$ where $p$ collects the
weight-exactly-$c$ factors and $q$ the weight-$\ge c+1$ factors (this respects
\eqref{eq:cf} because the factors are ordered by increasing index, hence by
non-decreasing weight). By the ($\Leftarrow$) direction $q\in\gamma_{c+1}(F)$, so
$p=gq^{-1}\in\gamma_{c+1}(F)$. But $p=\prod_{b\in B_c}b^{a_b}$ has image $\sum_b a_b[b]$
in $L_c=\gamma_c(F)/\gamma_{c+1}(F)$, which vanishes because $p\in\gamma_{c+1}(F)$. By
Proposition~\ref{prop:hall}, $L_c$ is free abelian on $\{[b]:b\in B_c\}$, so every
$a_b=0$: no weight-$c$ factor occurs, and every factor of $g$ has weight $\ge c+1$.
\end{proof}

\subsection*{The deduction}

\begin{proof}[Proof of Theorem~\ref{thm:main}]
Fix $c\ge1$. By Lemma~\ref{lem:easy}, $N_c\subseteq\gamma_c(F)$. For the reverse
inclusion, part (PR) of Theorem~\ref{thm:collection} states that the kernel of the
canonical surjection $F\to F/\gamma_c(F)$ is the normal closure
$\langle B_c\rangle^{F}=N_c$; that kernel is, by definition of the quotient,
$\gamma_c(F)$. Hence $\gamma_c(F)=N_c$. As $c$ was arbitrary, this is
Theorem~\ref{thm:main}; equivalently, by Proposition~\ref{prop:chain}, the single-level
statement \eqref{eq:S} holds.
\end{proof}

\begin{remark}[Logical status; non-circularity]\label{rem:noncirc}
By Proposition~\ref{prop:chain}, part (PR) of Theorem~\ref{thm:collection} is logically
\emph{equivalent} to Theorem~\ref{thm:main}; the substance of the present theorem resides
in this classical input, and we do not claim a derivation of it from anything strictly
weaker (Proposition~\ref{prop:irreducible} shows none can exist). What makes its use a
legitimate citation rather than an assumption of the conclusion is that
Theorem~\ref{thm:collection} is proved by a route logically prior to
Theorem~\ref{thm:main} that never mentions the subgroups $N_c$:
\begin{itemize}
\item \emph{Existence} of the collected form \eqref{eq:cf} and the finite-support clause
(C1) are the collection process: an algorithm rewriting a word in $X^{\pm}$ using only
the free-group identities \eqref{eq:id1}--\eqref{eq:id6}. On a length-$\ell$ word only
basic commutators of weight $\le\ell$ ever occur (finitely many), and each move strictly
decreases a well-founded complexity (the number of uncollected occurrences of
lower-order generators), so the process terminates. The residual high-weight tail is
annihilated by Magnus residual nilpotence $\bigcap_m\gamma_m(F)=1$ (W.~Magnus 1935;
M.~Hall, Thm.~11.2.5), proved by the power-series embedding
$x_i\mapsto1+X_i$ into $\mathbb Z\langle\langle X_1,\dots,X_r\rangle\rangle$.
\item \emph{Uniqueness} in (NF) is the $\mathbb Z$-linear independence of the leading
Magnus monomials of the basic commutators (equivalently the free-abelian layer clause of
Proposition~\ref{prop:hall}).
\item \emph{Presentation} (PR) follows from (NF) and (C1) by the standard argument,
which our Lemma~\ref{lem:support} renders transparent: by Lemma~\ref{lem:support} the
elements killed in $F/\gamma_c(F)$ are exactly those whose collected support lies in
weights $\ge c$, and (C1) guarantees this support is finite; the relations
$\{b=1:b\in B_c\}$ collapse precisely these elements --- the collection corrections being
basic commutators of strictly larger weight --- so the kernel is exactly $N_c$. The
finite-support clause (C1) is what lets this terminate inside every $F/\gamma_m(F)$
(as in Proposition~\ref{prop:finiteS}) and pass to $F$ via Magnus residual nilpotence.
\end{itemize}
None of these inputs mentions $N_c$. By contrast Proposition~\ref{prop:irreducible}
shows that no purely graded substitute can exist: the collapse is invisible to every
layer $\gamma_w/\gamma_{w+1}$, to the associated graded ring, and to the Magnus
embedding. The exact non-graded ingredient that breaks this symmetry is precisely the
finite-support clause (C1). The elementary content of
\S\S\ref{sec:easy}--\ref{sec:resnilp} pins down exactly what (C1)/(PR) must supply: the
elementary descent yields only $\gamma_c(F)\subseteq\bigcap_m N_c\gamma_m(F)$
(Lemma~\ref{lem:moddescent}, Proposition~\ref{prop:finiteS}), and the missing collapse
$\bigcap_m N_c\gamma_m(F)=N_c$ is exactly the statement that an element of $\gamma_c(F)$
has a \emph{finite} collected support concentrated in weights $\ge c$ --- which (C1)
provides.
\end{remark}



\begin{remark}[How the issue's worry is resolved]\label{rem:issue}
A weight-$w$ basic commutator $[u,v]$ with $w>c$ need not be a visible commutator of a
weight-$c$ element; the first such case is $[[x_2,x_1],[x_3,x_1]]\in B_4$ relative to
$c=3$ (Example~\ref{ex:worked}). The naive route rewrites $[u,v]$ modulo
$\gamma_{w+1}$ as an integral combination of weight-$w$ basic commutators and
iterates, giving only $\gamma_c\subseteq N_c\gamma_{c+m}$ for all $m$
(Lemma~\ref{lem:moddescent}, Proposition~\ref{prop:finiteS}); passing to the limit
would require $\bigcap_m N_c\gamma_{c+m}=N_c$, i.e.\ residual nilpotence of $F/N_c$,
which is false for general normal subgroups (Example~\ref{ex:Z}) and not implied by
residual nilpotence of $F$ alone (Proposition~\ref{prop:irreducible}). Our treatment:
\begin{itemize}
\item proves the weight statement \emph{completely inside every nilpotent quotient}
$F/\gamma_m$ by a terminating downward induction on weight, including the issue's
example (Proposition~\ref{prop:finiteS}, Example~\ref{ex:worked});
\item isolates the remaining content as the collapse of the intersection,
equivalently residual nilpotence of $F/N_c$ (Proposition~\ref{prop:equivresnilp});
\item proves that this collapse cannot be obtained from graded data
(Proposition~\ref{prop:irreducible});
\item discharges it by the named, hypotheses-verified classical Hall collection process
(Theorem~\ref{thm:collection}, part (PR)), whose collection-process proof supplies the
exact non-graded ingredient (finite support of the collected coordinate vector) and is
established independently of Theorem~\ref{thm:main} (Remark~\ref{rem:noncirc}).
\end{itemize}
\end{remark}

\end{document}
